% Generated draft: requires mathematical and source review.
\documentclass[11pt]{article}
\usepackage[margin=25mm]{geometry}
\usepackage{fontspec}
\setmainfont{Latin Modern Roman}
\usepackage{amsmath,amssymb,mathtools,unicode-math}
\setmathfont{Latin Modern Math}
\usepackage{enumitem,needspace,xurl,hyperref}
\setlength{\parindent}{0pt}
\setlength{\parskip}{5pt}
\setlength{\emergencystretch}{3em}
\newcommand{\N}{\mathbb{N}}
\newcommand{\Z}{\mathbb{Z}}
\newcommand{\Q}{\mathbb{Q}}
\newcommand{\R}{\mathbb{R}}
\newcommand{\C}{\mathbb{C}}
\newcommand{\F}{\mathbb{F}}
\newcommand{\A}{\mathbb{A}}
\newcommand{\PP}{\mathbb{P}}
\newcommand{\E}{\mathbb{E}}
\newcommand{\eps}{\varepsilon}
\newcommand{\dd}{\,\mathrm{d}}
\begin{document}
\section*{UnsolvedMath ID: 700550}
\textbf{Source problem code:} catalogue entry
\subsection{Definitions and mathematical statement}
Let \(G\) be a finite Chevalley group. Relative to its standard \(BN\)-pair, with Borel subgroup \(B\), a subgroup \(P\leq G\) is parabolic if it contains a conjugate of \(B\). For subgroups \(P,L\leq G\), define their setwise product by
\[
PL=\{p\ell:p\in P,\ \ell\in L\}.
\]
The requested classification is, for every such \(G\), to describe all members of
\[
\mathcal{L}(G)=\{L\leq G:\text{there exists a parabolic subgroup }P\leq G\text{ with }G=PL\},
\]
using arguments that do not invoke the classification of finite simple groups (CFSG). Equivalently, for such a \(P\), the right multiplication action of \(L\) on the right coset space \(P\backslash G\), given by \((Pg)\cdot\ell=Pg\ell\), must be transitive.

\textbf{Unresolved qualifications.} Under the definition above, \(G\) itself is parabolic. If \(P=G\) is allowed, every subgroup \(L\leq G\) satisfies the displayed condition. The source does not explicitly require \(P<G\); whether proper parabolic subgroups are intended requires clarification. The source also does not specify the precise convention for finite Chevalley groups, including whether twisted groups are included or which central forms and quotients are allowed. These choices remain unresolved, and no restriction to a particular convention is imposed here.
\subsection{Short English statement}
For each finite Chevalley group \(G\), classify the subgroups \(L\) for which \(G=PL\) for some parabolic subgroup \(P\), without using CFSG. The intended scope of Chevalley groups and whether \(P\) must be proper require clarification; allowing \(P=G\) makes every subgroup eligible.
\subsection{Sources}
\begin{itemize}
\item Local UnsolvedMath registry: supplied complete source record, permanent ID \(700550\), exact source problem code \texttt{catalogue entry}, title \emph{Kourovka Notebook Problem 9.15}; upstream revision \texttt{d371ee917ac9647fbb06cd6fc529deec6f3dda82}. The supplied registry labels the problem \texttt{open}; this is a supplied status, not an independently verified mathematical assessment.
\item \emph{The Kourovka Notebook (Unsolved Problems in Group Theory)}, original problem 9.15, 9th issue (1984); cited in the supplied record through the 21st issue edition (2026), \texttt{arXiv:1401.0300v48}. Proposed by A. S. Kondratiev.
\item The supplied collection description names E. I. Khukhro and V. D. Mazurov as editors of the Kourovka Notebook collection covering Issues 1--21 (1965--2026). The cited source and edition have not been independently consulted for this draft.
\end{itemize}
\end{document}
